\chapter{相关研究基础}
\label{chap:foundation}

本章从建筑幕墙巡检与数字化运维、幕墙缺陷检测技术、AI Agent 工程范式和韧性评估报告生成四个方向梳理本系统的相关研究基础。现有方法在单类缺陷检测能力上已经具有一定积累，但在项目级任务编排、Agent 自主任务调度和 LLM 分层信息整合方面仍存在一定空白。

\section{建筑幕墙巡检与数字化运维}
\label{sec:foundation-building-curtain-wall-inspection-and-digital-operation-and-maintenance}

建筑幕墙巡检长期以来依赖人工现场勘查。人工巡检能够结合经验进行综合判断，但在高层建筑、大面积幕墙和复杂环境下效率较低，且存在安全风险。随着无人机技术和图像采集设备的发展，基于无人机的建筑立面巡检逐渐成为重要方向。Freimuth 和 K{\"o}nig 对无人机建筑检查的规划和执行流程进行了系统研究，说明无人机能够降低高空检测成本并提升数据采集效率\cite{freimuth2018uav}。在此基础上，红外热成像、图像配准和 BIM、AI 结合等方法进一步扩展了建筑立面检测的数据来源、空间定位能力和智能分析能力\cite{li2024facade,zhang2023uavbim,leng2018bimai}。

但是，数字化采集并不等于智能化评估。无人机可以快速采集图像，但图像采集之后仍需要图像管理、缺陷检测、结果解释和报告生成。若缺少后续软件系统，采集到的图像仍然是离散资料，难以转化为运维决策。因此，建筑幕墙评估系统需要将采集后的图像组织为项目资产，并建立从图像收集到评估报告生成的完整处理流程。

从工程应用角度看，采集设备能够提高图像获取效率，检测模型能够提高局部缺陷识别效率，但真正面运维单位和工程师的交付物通常是带有结论、依据和建议的评估报告。这意味着系统需要同时处理视觉证据、结构化指标、项目语义和文本表达。若系统只停留在缺陷框或掩码图层面，仍然需要人工把这些结果重新解释、筛选、汇总和撰写。因此，图像检测之后的组织、解释和报告生成，是建筑幕墙从数字化运维走向智能化运维的关键环节。

幕墙工程对象同时具有明显的项目属性。同一种缺陷出现在不同区域，其工程意义可能不同；同一张图像中的不同材料，也可能对应不同检测任务。传统视觉检测研究往往把图像视作独立样本，而实际运维更关心项目内图像之间的关系、图像组与建筑区域的关系、缺陷分布与维护优先级之间的关系等。本文的系统设计正是围绕这种项目级应用需求展开。

\section{幕墙缺陷检测技术}
\label{sec:foundation-curtain-wall-defect-detection-technology}

在本文讨论的幕墙巡检场景中，图像是最直接也最容易获得的数据来源。算法面对的通常不是标准实验图，而是带有反光、阴影、倾斜视角、材料纹理和背景遮挡的现场照片。因此，缺陷检测一般要完成几件比较实际的工作：判断缺陷类型，给出缺陷位置，并尽量形成面积、长度、平整度或置信度等可解释指标。早期或轻量化场景中仍会使用传统计算机视觉方法，例如利用 Canny 边缘检测和 Hough 直线检测提取轮廓或线性结构，也可以结合 Sobel 梯度、Laplacian 方差、傅里叶频域分析和形态学开闭运算描述纹理突变、模糊程度或局部形态变化\cite{canny1986edge,duda1972hough}。这类方法的优点是实现过程清楚，对大规模标注数据依赖较小，在裂缝、边缘和规则构件分析中也比较容易解释；不足是对现场条件较敏感，换到另一种光照、另一类材质或另一批拍摄角度后，阈值和规则往往需要重新调整。

深度学习方法把特征提取更多交给数据本身。卷积神经网络不再手工指定每一个边缘或纹理规则，而是在训练过程中逐层形成从低级纹理到高层语义的图像表征；ResNet 等残差网络进一步通过残差连接缓解深层网络训练困难，因而常被用于图像分类任务\cite{article1,conf1}。如果任务关注“缺陷在哪里”，目标检测模型会更直接，它可以同时给出类别和位置。基于候选区域的两阶段方法，以及以 YOLO、SSD 为代表的一阶段方法，是目标检测中较常见的两条路线；国内已有综述对 R-CNN、Faster R-CNN、YOLO 和 SSD 等代表性算法进行了梳理\cite{fan2020objectdetection}。其中，YOLO 系列强调端到端和实时检测，在需要快速定位缺陷候选区域时较为适用\cite{redmon2016yolo}。

当评估需要更细的边界或面积信息时，仅有检测框往往不够，语义分割和实例分割模型会更有价值。U-Net 等编码器-解码器结构是像素级分割中的典型思路，相关文献综述已经对深度学习语义分割中的区域分类、像素分类、全监督和弱监督方法进行了总结\cite{tian2019semanticsegmentation}。DeepLabV3+ 通过空洞卷积和编码器-解码器结构提升语义分割效果\cite{chen2018deeplabv3plus}；视觉 Transformer 借助自注意力机制增强长距离依赖建模能力，为视觉识别、目标检测和视觉分割提供了新的表征基础\cite{zhou2023visiontransformer}；SegFormer 则进一步将 Transformer 表征用于高效语义分割\cite{xie2021segformer}。此外，Mask R-CNN 等实例分割方法和 Segment Anything 分别从实例掩码预测、通用分割能力等方向扩展了区域分割的使用范围\cite{su2022instancesegmentation,kirillov2023sam}。放到建筑结构和幕墙巡检中，裂缝检测、玻璃检测和表面缺陷识别已经有较多积累。Ye 等人利用深度学习方法进行结构裂缝检测，可为石材或混凝土裂缝识别提供参考\cite{ye2019crack}；Mei 等人研究真实场景中的玻璃检测问题，对玻璃幕墙区域分析也具有参考意义\cite{mei2020glass}。

\begin{figure}[!htb]
  \centering
  \resizebox{\linewidth}{!}{%
    \begin{tikzpicture}[
      font=\scriptsize,
      >=stealth,
      lane/.style={draw, rounded corners=2pt, line width=0.45pt, fill=gray!4},
      capability/.style={draw, rounded corners=2pt, align=center, minimum width=2.2cm, minimum height=0.58cm, line width=0.45pt, fill=blue!6, inner sep=2pt},
      method/.style={draw, rounded corners=2pt, align=center, minimum width=3.35cm, minimum height=0.58cm, line width=0.45pt, fill=orange!10, inner sep=2pt},
      output/.style={draw, rounded corners=2pt, align=center, minimum width=3.3cm, minimum height=0.58cm, line width=0.45pt, fill=green!8, inner sep=2pt},
      arrow/.style={->, line width=0.45pt},
      header/.style={font=\scriptsize\bfseries, align=center}
    ]
      \node[header] at (1.1, 0) {原子检测能力};
      \node[header] at (4.7, 0) {核心技术路线};
      \node[header] at (8.85, 0) {结构化输出};

      \draw[lane] (-0.25, -0.38) rectangle (10.75, -1.32);
      \node[capability] (corrosion) at (1.1, -0.85) {金属锈蚀};
      \node[method] (corrosion-method) at (4.7, -0.85) {YOLO / DeepLabV3+\\检测与语义分割};
      \node[output] (corrosion-output) at (8.85, -0.85) {锈蚀区域、置信度\\锈蚀像素数与占比};

      \draw[lane] (-0.25, -1.46) rectangle (10.75, -2.40);
      \node[capability] (crack) at (1.1, -1.93) {石材裂缝};
      \node[method] (crack-method) at (4.7, -1.93) {YOLO / SAM / SegFormer\\候选区域与裂缝分割};
      \node[output] (crack-output) at (8.85, -1.93) {裂缝掩码、覆盖图\\裂缝像素数与占比};

      \draw[lane] (-0.25, -2.54) rectangle (10.75, -3.48);
      \node[capability] (stain) at (1.1, -3.01) {石材污渍};
      \node[method] (stain-method) at (4.7, -3.01) {YOLO 实例分割\\透视校正与局部分析};
      \node[output] (stain-output) at (8.85, -3.01) {污渍区域、校正图\\平均与最高污渍占比};

      \draw[lane] (-0.25, -3.62) rectangle (10.75, -4.56);
      \node[capability] (flatness) at (1.1, -4.09) {玻璃平整度};
      \node[method] (flatness-method) at (4.7, -4.09) {Canny / Hough / Sobel / FFT\\多特征规则融合};
      \node[output] (flatness-output) at (8.85, -4.09) {不平整判定指标\\边缘 / 直线 / 梯度 / 频谱};

      \draw[lane] (-0.25, -4.70) rectangle (10.75, -5.64);
      \node[capability] (spalling) at (1.1, -5.17) {玻璃爆裂};
      \node[method] (spalling-method) at (4.7, -5.17) {ResNet34 二分类\\数据清洗与增强训练};
      \node[output] (spalling-output) at (8.85, -5.17) {爆裂 / 未爆裂判定\\分类置信度};

      \draw[arrow] (corrosion) -- (corrosion-method);
      \draw[arrow] (corrosion-method) -- (corrosion-output);
      \draw[arrow] (crack) -- (crack-method);
      \draw[arrow] (crack-method) -- (crack-output);
      \draw[arrow] (stain) -- (stain-method);
      \draw[arrow] (stain-method) -- (stain-output);
      \draw[arrow] (flatness) -- (flatness-method);
      \draw[arrow] (flatness-method) -- (flatness-output);
      \draw[arrow] (spalling) -- (spalling-method);
      \draw[arrow] (spalling-method) -- (spalling-output);
    \end{tikzpicture}%
  }
  \caption{五类原子检测能力核心技术路线总览}
  \label{fig:foundation-atomic-detection-capability-technical-routes}
\end{figure}

如\cref{fig:foundation-atomic-detection-capability-technical-routes} 所示，Agent 自主调度的五类原子检测能力来自团队前期积累的检测模型和实验脚本，本系统将这些模型抽象为可调度、可替换、可治理的原子检测能力。在系统链路中，原子检测能力负责给出局部视觉判断，Agent 自主任务调度方法负责选择应执行的能力，LLM 分层信息整合策略负责将局部结果组织为图像总结和项目报告。本节对五类原子检测能力的基本原理进行说明。

\subsection{金属锈蚀检测原理}
\label{sec:foundation-metal-corrosion-detection-principle}

金属构件发生锈蚀以后，在图像上一般会表现为颜色变深、边界不规则、局部纹理粗糙等现象。系统对这一类问题要做的事情比较明确：把锈蚀区域找出来，同时给出锈蚀区域数量、置信度、锈蚀像素数和锈蚀占比等指标。为了兼顾检测速度和区域精度，当前金属锈蚀检测没有只固定在一种模型上，而是保留了两条路径。一条是基于 YOLO 的检测或分割路径，比较适合快速找出锈蚀候选区域，并输出缺陷框、掩码和类别置信度；另一条是基于 DeepLabV3+ 的语义分割路径，它把锈蚀检测看作二分类像素分割问题，通过编码器提取图像特征，再借助空洞卷积扩大感受野，去识别形状并不规则的锈蚀区域\cite{chen2018deeplabv3plus}。在代码实现里，DeepLabV3+ 使用 MobileNetV2 作为轻量化骨干网络，这样能把推理成本压低一些，也更便于后续工程部署\cite{sandler2018mobilenetv2}。

实际推理时，脚本会先读取原始图像，再完成尺度变换、张量化和归一化处理，之后得到锈蚀概率图。走语义分割路径时，二值掩码不是简单固定阈值得到的，而是会根据概率分布来判断：整体概率较高时使用 Otsu 自适应阈值，其他情况下使用置信度阈值对概率图进行二值化。为了减少零散噪声，脚本还会做形态学开运算，并根据连通域或轮廓面积把过小区域过滤掉。经过这些处理后，系统可以从掩码里统计锈蚀像素数、锈蚀面积占比、区域数量和平均置信度，同时生成锈蚀覆盖图给前端展示。这样一来，锈蚀检测结果既有图像证据，也有可以继续参与项目级评估的结构化指标。

\subsection{石材裂缝检测原理}
\label{sec:foundation-stone-crack-detection-principle}

石材或混凝土幕墙上的裂缝，在现场图像中并不总是很明显。它可能很细，也可能被石材本身的纹理、阴影、接缝或者拍摄噪声遮住一部分，所以只靠边缘检测或固定阈值很容易把纹理误认为裂缝，或者漏掉对比度较低的真实裂缝。基于这个原因，裂缝检测脚本采用了深度分割路线，把裂缝识别处理成像素级语义分割任务。原始系统里并不是直接对整张图做一次判断，而是包含幕墙区域检测、区域分割和裂缝分割等步骤；其中可以结合 YOLO 旋转框模型、MobileSAM 或 SegFormer 等能力，先得到相对稳定的石材候选区域，再交给裂缝分割模型输出裂缝掩码\cite{xie2021segformer,kirillov2023sam}。

进入裂缝分割阶段后，图像会经过预处理并输入分割模型，模型给出的结果是每个像素属于裂缝类别的概率。后处理部分会把概率图转成二值掩码，再完成尺寸恢复和高亮叠加，形成裂缝覆盖图、裂缝掩码图和裂缝区域统计结果。与只输出目标检测框相比，像素级掩码能把裂缝的长度、宽度和空间走向表达得更细一些，也能进一步计算裂缝像素数和裂缝占比。对于后续评估来说，这种结果比“某处有裂缝”的简单判断更有用。

\subsection{石材污渍检测原理}
\label{sec:foundation-stone-stain-detection-principle}

石材污渍和裂缝的表现方式不太一样。裂缝更像细长线状结构，而污渍通常是污染、渗水痕迹或颜色异常造成的块状、片状变化，边界不一定清楚，还会受到拍摄角度和透视变形的影响。如果直接在整张图上计算污渍比例，不同幕墙块之间的倾斜程度会把结果带偏。因此，污渍检测脚本采用了“幕墙块分割、透视校正、局部污渍分析”的组合路线。具体做法是，先用 YOLO 模型对图像中的幕墙块进行实例分割，得到每个候选幕墙块的掩码；再根据掩码轮廓计算凸包，并通过多边形近似寻找四个角点；角点按照左上、右上、右下、左下的顺序排好以后，系统就能利用透视变换把倾斜拍摄的幕墙块校正成正视局部图像。

经过透视校正后，每个局部幕墙块会被送入传统图像处理算法，用来估计污渍面积。原始脚本会调用污渍分析模块，计算每个校正块的污渍比例，并生成对应的结果图像。这里输出的图像也有不同用途：整体标注图用来查看所有候选幕墙块的位置，局部校正图用来确认检测区域，结果图则展示污渍分析后的效果。这样处理的好处在于，深度学习模型主要负责区域定位，传统图像处理负责比例计算，二者分工比较清楚。由此得到的“平均污渍占比”和“最高污渍占比”，也就不只是整图上的粗略统计，而是有了更明确的区域含义。

\subsection{玻璃平整度检测原理}
\label{sec:foundation-glass-flatness-detection-principle}

玻璃平整度检测主要关注玻璃幕墙中可能出现的局部变形、反射异常或不平整现象。这类问题和裂缝、锈蚀不同，图像上往往没有一个清楚的缺陷边界。由于玻璃本身透明、反光、纹理少，真实的不平整更可能体现为边缘畸变、直线偏折、梯度分布异常，或者频域能量变化。基于这些特点，脚本没有直接采用单一分类判断，而是采用“玻璃区域识别 + 多特征规则判别”的实现路径。也就是说，系统会先通过玻璃检测或分割模型得到玻璃区域，再把后续分析集中在局部区域内，减少背景和窗框带来的干扰。

在局部玻璃区域中，脚本会先进行裁剪，然后使用 Canny 方法提取边缘，并计算边缘数量和 Laplacian 方差，用来反映局部纹理和清晰度变化\cite{canny1986edge}。直线分析这一部分会通过概率 Hough 变换检测线段，再统计线段角度标准差，以判断反射线或边缘线有没有出现异常偏折\cite{duda1972hough}。梯度分析则利用 Sobel 算子计算梯度幅值，并统计梯度均值和梯度标准差；频域分析会对灰度图执行离散傅里叶变换，观察频谱能量差异。系统把直线、梯度和频域三个判别结果放在一起，通过多数投票得到玻璃是否不平整的结论。这样做虽然不如端到端模型简洁，但每一个判断依据都比较容易解释。

\subsection{玻璃爆裂检测原理}
\label{sec:foundation-glass-spalling-detection-principle}

玻璃爆裂检测在系统接口上看起来比较简单，本质上是判断玻璃幕墙图像中是否存在破裂、爆裂或碎裂风险。脚本采用的是基于 ResNet34 的二分类模型，把输入图像划分为“存在缺陷”和“不存在缺陷”两类。模型以 ResNet34 作为骨干网络，并把最后的全连接层替换成二分类输出层。推理时，图像会被缩放到 $224\times224$，再转换为张量，并按照 ImageNet 均值和标准差进行归一化，之后输入模型得到类别预测。

不过，这项能力并不只是一次分类推理。为了让二分类模型在真实拍摄条件下更稳定，脚本前面还包含了数据准备和训练前处理流程。数据清洗阶段会检查图像完整性、文件重复、分辨率、标注缺失和拍摄角度；数据增强阶段会使用旋转、水平翻转、缩放、亮度对比度扰动、高斯噪声、椒盐噪声和透视畸变等方式，让模型见过更多可能出现的图像变化；数据划分阶段则按照工程类型、区域和年份构造分层键，再把数据划分为训练集、验证集和测试集。原始脚本中还保留了基于 Canny 与 Hough 直线的辅助处理，用来过滤与横纵轴近似平行的线性纹理干扰。因此，玻璃爆裂检测虽然对上层系统表现为一个二分类接口，但它背后仍然依赖数据治理、数据增强和模型训练等一整套流程。

综上，五类原子检测能力在技术路线和输出形态上存在差异：金属锈蚀和石材裂缝更偏向分割，石材污渍结合了实例分割与透视校正，玻璃平整度依赖多种视觉特征融合，玻璃爆裂则采用深度分类模型。本文将这些差异封装在统一的原子能力接口之后，使上层 Agent 不需要感知每个算法的内部实现，而只需要基于图像内容选择任务代码并接收结构化结果。这也是本文从单点检测算法走向工程化智能评估系统的基础。

\section{从单点检测到项目级评估的差距}
\label{sec:foundation-single-algorithm-to-project-level-assessment-gap}

已有幕墙缺陷检测研究通常以单一任务为目标，例如裂缝识别、玻璃检测或表面污渍检测。这类研究能够回答算法层面的问题：给定输入图像，模型是否能够准确输出目标区域。然而，项目级评估需要回答更复杂的问题：哪些图像需要执行哪些检测，多个检测结果之间如何互相补充，哪些缺陷对项目整体风险影响更大，报告中应如何表达证据与结论之间的关系。

\begin{figure}[!htb]
  \centering
  \resizebox{\linewidth}{!}{%
    \begin{tikzpicture}[
      font=\scriptsize,
      >=stealth,
      outer/.style={draw, rounded corners=2pt, line width=0.45pt, fill=gray!4},
      single/.style={draw, rounded corners=2pt, align=center, minimum width=2.75cm, minimum height=0.66cm, line width=0.45pt, fill=blue!6, inner sep=2pt},
      gap/.style={draw, rounded corners=2pt, align=center, minimum width=4.0cm, minimum height=0.62cm, line width=0.45pt, fill=orange!10, inner sep=2pt},
      target/.style={draw, rounded corners=2pt, align=center, minimum width=2.95cm, minimum height=0.62cm, line width=0.45pt, fill=green!8, inner sep=2pt},
      title/.style={font=\scriptsize\bfseries, align=center},
      arrow/.style={->, line width=0.45pt}
    ]
      \draw[outer] (-0.15, 0.45) rectangle (12.25, -4.25);

      \node[title] at (1.65, 0) {单点检测算法};
      \node[title] at (6.25, 0) {必须跨越的缺口};
      \node[title] at (10.35, 0) {项目级评估问题};

      \node[single] (single-image) at (1.65, -1.00) {输入对象\\单张图像};
      \node[single] (single-result) at (1.65, -2.18) {输出结果\\缺陷框 / 掩码 / 分类};
      \node[single, minimum height=0.7cm] (single-question) at (1.65, -3.28) {核心回答\\这张图有什么缺陷};
      \draw[arrow] (single-image) -- (single-result);
      \draw[arrow] (single-result) -- (single-question);

      \node[gap] (gap-data) at (6.25, -0.80) {数据组织缺口\\图像进入项目 / 图像组 / 资产};
      \node[gap] (gap-task) at (6.25, -1.78) {任务编排缺口\\材质识别与检测能力选择};
      \node[gap] (gap-explain) at (6.25, -2.68) {结果解释缺口\\像素坐标指标需转化为工程语义};
      \node[gap] (gap-report) at (6.25, -3.58) {报告生成缺口\\局部结论需融合为建筑级报告};

      \node[target] (target-select) at (10.35, -0.80) {哪些图像需要检测};
      \node[target] (target-distribution) at (10.35, -1.78) {缺陷如何分布};
      \node[target] (target-risk) at (10.35, -2.68) {风险如何判断};
      \node[target] (target-report) at (10.35, -3.58) {报告如何生成};

      \coordinate (gap-branch) at (3.64, -3.28);
      \draw[arrow] (single-question.east) -- (gap-branch) |- (gap-data.west);
      \draw[arrow] (single-question.east) -- (gap-branch) |- (gap-task.west);
      \draw[arrow] (single-question.east) -- (gap-branch) |- (gap-explain.west);
      \draw[arrow] (single-question.east) -- (gap-branch) |- (gap-report.west);

      \draw[arrow] (gap-data) -- (target-select);
      \draw[arrow] (gap-task) -- (target-distribution);
      \draw[arrow] (gap-explain) -- (target-risk);
      \draw[arrow] (gap-report) -- (target-report);
    \end{tikzpicture}%
  }
  \caption{从单点检测到项目级评估的跨越关系}
  \label{fig:foundation-single-detection-to-project-assessment-transition}
\end{figure}

如\cref{fig:foundation-single-detection-to-project-assessment-transition} 所示，从单点检测到项目级评估之间至少存在四个缺口。第一是数据组织缺口，图像需要按照项目、图像组、图像和检测任务多层组织。第二是任务编排缺口，系统需要根据图像内容动态选择检测能力，而不是机械执行全部模型。第三是结果解释缺口，模型输出的像素和坐标指标需要转化为工程语言。第四是报告生成缺口，系统需要把多张图像的局部结论融合为项目级评估意见。本文的研究价值正是在这些缺口处构建桥梁。

如果只关注检测模型，系统会缺少业务闭环；如果只关注 Web 管理界面，系统又缺少智能决策能力。本文选择的路线是以工程系统组织检测模型，以 Agent 承担语义决策和信息整合，从而把二者连接为完整应用。

\section{AI Agent 的发展与工程范式}
\label{sec:foundation-artificial-intelligence-agent-development-and-engineering-paradigm}

大语言模型的兴起使 Agent 成为当前人工智能应用的重要范式。早期大语言模型主要用于文本生成和问答，而 Agent 更强调模型与外部工具、环境和用户目标之间的交互。ReAct 提出将推理和行动交替结合，使模型能够生成思考过程并执行任务动作\cite{yao2023react}。Toolformer 进一步说明语言模型可以学习何时调用工具、如何传递参数以及如何利用工具返回结果\cite{schick2023toolformer}。这些工作共同说明，模型本身的语言能力需要通过工具调用机制转化为可执行能力。

工业界和开源社区也逐渐形成 Agent 工作流开发范式。AutoGen 通过多智能体对话组织复杂任务，说明 Agent 可以被看作具有不同角色和能力的可编程组件\cite{wu2023autogen}。相关综述指出，大语言模型 Agent 可应用于单智能体、多智能体和人机协作等场景，其关键能力包括规划、记忆、工具使用和反馈学习\cite{wang2023llmagents}。在实践中，Agent 不会完全自由行动，而是被放置在受控工作流中，通过结构化输入输出、工具白名单和状态管理保证系统可控。

本文正是基于这一认识，将 Agent 设计为业务链路中的受控智能节点。分类 Agent 不是自由选择任意工具，而是在预定义五类任务代码中选择；总结 Agent 不是直接修改系统状态，而是通过回调接口提交总结结果；项目报告 Agent 不是自行读取数据库，而是接收由 BIZ 服务组织好的项目上下文。这样的设计符合当前 Agent 工程化趋势：让模型承担语义理解和决策任务，让工程系统负责边界和状态。

\begin{table}[!htb]
  \caption{大语言模型应用范式的演进}\label{tab:llm-agent-evolution}
  \centering
  \begin{tabular}{@{}p{0.12\linewidth}p{0.4\linewidth}p{0.4\linewidth}@{}} \toprule
    \textbf{范式} & \textbf{主要特征} & \textbf{局限与本文启发} \\ \midrule
    文本生成 & 模型根据提示词直接输出自然语言 & 适合问答和撰写，但难以与业务状态和外部工具形成闭环 \\
    检索增强 & 先检索外部知识，再生成回答 & 能补充上下文，但仍偏向知识问答，不直接解决任务调度问题 \\
    工具调用 & 模型根据目标选择外部 API 或工具 & 适合连接模型与工程能力，但需要严格约束工具集合和参数 \\
    工作流 Agent & Agent 被放入确定业务流程，承担局部决策 & 更适合工程落地，本文采用该路线实现检测路由和信息整合 \\ \bottomrule
  \end{tabular}
\end{table}

\cref{tab:llm-agent-evolution} 展示了大语言模型应用范式从文本生成到工作流 Agent 的变化。本文将 Agent 视为工程工作流中的智能决策节点，而非完全自主的软件机器人。这样的设计更符合垂直行业系统的要求：业务流程必须可解释，数据状态必须可追踪，外部工具调用必须可控，失败情况必须可恢复。

\section{Agent 工程化的可控性问题}
\label{sec:foundation-agent-engineering-controllability-problem}

Agent 系统虽然具备更强的自主性，但也带来工程可控性问题。大语言模型输出具有概率性，可能出现格式不稳定、内容越界、幻觉或工具选择错误等情况。幻觉问题已经被大量研究指出，其核心风险在于模型可能生成看似合理但缺少事实依据的内容，进而影响真实系统中的信息可靠性\cite{huang2023hallucination}。在工具调用场景中，模型还需要判断何时调用工具、调用哪个工具以及如何组织参数。Toolformer、ToolLLM 等研究说明了大语言模型可以学习或增强工具使用能力，但也从侧面说明工具使用不是天然稳定能力，需要通过工具集合、调用路径和评估机制进行约束\cite{schick2023toolformer,qin2023toolllm}。对于普通问答系统，这些问题可能只是回答质量问题；对于工程业务系统，如果 Agent 输出直接驱动数据库状态或任务调度，则可能造成错误的业务流程。因此，Agent 工程化必须关注结构化输出、工具白名单、状态隔离和人工复核。

\begin{figure}[!htb]
  \centering
  \resizebox{\linewidth}{!}{%
    \begin{tikzpicture}[
      font=\scriptsize,
      >=stealth,
      risk/.style={draw, rounded corners=2pt, align=center, minimum width=4.0cm, minimum height=0.58cm, line width=0.45pt, fill=red!6, inner sep=2pt},
      control/.style={draw, rounded corners=2pt, align=center, minimum width=5.1cm, minimum height=0.58cm, line width=0.45pt, fill=green!8, inner sep=2pt},
      group/.style={draw, rounded corners=2pt, line width=0.45pt, fill=gray!4},
      layer/.style={font=\scriptsize\bfseries, align=center},
      principle/.style={font=\scriptsize, align=center},
      arrow/.style={->, line width=0.5pt}
    ]
      \path[use as bounding box] (-0.60, 1.05) rectangle (12.30, -5.02);
      \node[principle] at (5.85, 0.78) {核心原则：Agent 承担受控决策，工程系统负责边界、状态与恢复};

      \draw[group] (0.00, 0.48) rectangle (4.90, -4.92);
      \draw[group] (5.70, 0.48) rectangle (11.70, -4.92);

      \node[layer] at (2.45, 0.08) {Agent 工程化风险};
      \node[layer] at (8.70, 0.08) {工程约束机制};

      \node[risk] (risk-format) at (2.45, -0.60) {输出格式不稳定};
      \node[control] (control-format) at (8.70, -0.60) {JSON Schema / 固定字段 / 解析校验};
      \draw[arrow] (risk-format) -- (control-format);

      \node[risk] (risk-tool) at (2.45, -1.36) {工具选择错误};
      \node[control] (control-tool) at (8.70, -1.36) {工具白名单 / 有限任务代码};
      \draw[arrow] (risk-tool) -- (control-tool);

      \node[risk] (risk-hallucination) at (2.45, -2.12) {幻觉与事实偏移};
      \node[control] (control-hallucination) at (8.70, -2.12) {结构化上下文 / 系统组装数据源};
      \draw[arrow] (risk-hallucination) -- (control-hallucination);

      \node[risk] (risk-state) at (2.45, -2.88) {状态误推进};
      \node[control] (control-state) at (8.70, -2.88) {BIZ 统一推进状态 / Agent 不直接写库};
      \draw[arrow] (risk-state) -- (control-state);

      \node[risk] (risk-recovery) at (2.45, -3.64) {异常不可恢复};
      \node[control] (control-recovery) at (8.70, -3.64) {状态机 / 失败记录 / 重试机制};
      \draw[arrow] (risk-recovery) -- (control-recovery);

      \node[risk] (risk-trust) at (2.45, -4.40) {结果可信度不足};
      \node[control] (control-trust) at (8.70, -4.40) {人工复核 / 评论修正 / 审核闭环};
      \draw[arrow] (risk-trust) -- (control-trust);

    \end{tikzpicture}%
  }
  \caption{Agent 工程化风险与约束机制关系}
  \label{fig:foundation-agent-engineering-risk-control-mechanism}
\end{figure}

如\cref{fig:foundation-agent-engineering-risk-control-mechanism} 所示，现有研究和工程实践中的解决思路大体可以分为三类。第一类是限制动作空间，即只允许 Agent 在预定义工具、API 或任务代码中选择，ReAct 通过推理与行动交替使模型能够形成可解释行动轨迹，AutoGen 则通过多 Agent 对话和人类参与机制组织复杂任务流程\cite{yao2023react,wu2023autogen}。第二类是约束输出格式，通过语法、解析器或有限状态机限制生成结果，使模型输出满足结构化接口要求。PICARD 通过增量解析约束自回归解码，Outlines 相关研究则将文本生成表示为有限状态机转移，以正则表达式或上下文无关文法保证生成结构\cite{scholak2021picard,willard2023efficient}。第三类是将 Agent 放入受控工作流，由工程系统负责状态推进、失败重试、日志记录和人工复核。相关综述也指出，LLM Agent 的规划、工具使用和反馈机制需要与评估和控制机制共同设计，才能进入可靠应用\cite{wang2023llmagents}。

本文采用的思路是将 Agent 决策限制在有限动作空间内。分类 Agent 只能输出五种任务代码的组合；图像总结 Agent 的结果只作为总结字段保存，不直接修改检测状态；项目报告 Agent 接收的输入由 BIZ 服务生成，输出也通过固定回调进入数据库。这样做牺牲了一部分开放性，但换来了工程系统所需的确定性和可追踪性。对于建筑幕墙评估这种专业场景，这种受控形式是 Agent 自主任务调度方法能够落地的基础。

\section{工程软件中的异步任务编排}
\label{sec:foundation-engineering-software-asynchronous-task-orchestration}

长耗时任务编排是智能检测系统必须解决的问题。模型推理、外部 Agent 调用、对象上传和报告生成都可能持续数秒到数分钟，若采用同步 HTTP 请求等待全部结果，系统容易超时，用户体验也较差。成熟工程系统通常采用任务队列、Worker、回调和状态查询机制，将“启动任务”和“获得结果”解耦。

本文系统采用类似思想，但结合 Agent 场景进行了扩展。用户启动检测后，BIZ 服务只负责创建任务并入队；分类、推理和总结能力由不同 Activity 服务异步完成；完成后通过回调上报结果；前端通过 WebSocket 感知状态变化。该模式使 Agent 和模型调用可以作为后台任务运行，也使系统能够在失败时记录状态并支持重试。

异步任务编排对本系统尤为重要，因为 Agent 调用和模型推理都具有不可忽略的不确定性。Agent 平台可能因为网络、限流或输出异常导致失败；视觉模型可能因为图像格式、模型加载或运行环境导致失败；对象存储上传也可能受到网络影响。如果所有步骤都放在一次同步请求中，任何环节失败都会导致整个请求失败，且用户无法获得过程反馈。通过任务状态机和回调机制，系统能够把这些不确定因素转化为可记录、可展示和可恢复的状态变化。

从软件工程角度看，异步任务不仅是性能优化，也是复杂业务建模方式。检测主任务、检测子任务、总结任务和报告任务之间存在依赖关系。只有把这些关系显式建模为数据库状态和 Worker 编排逻辑，系统才能保证“分类失败不继续检测”“检测子任务全部成功才总结”“无检测子任务时直接成功”等业务规则被稳定执行。

\section{LLM 信息整合与报告生成}
\label{sec:foundation-large-language-model-information-integration-and-report-generation}

大语言模型擅长自然语言生成，但在知识密集型任务中，输入信息的组织方式直接影响输出质量。检索增强生成通过外部知识库向模型提供上下文，缓解模型仅依赖参数知识导致的不可靠问题\cite{lewis2020rag}。与此类似，工程报告生成也需要将外部数据组织成适合模型处理的上下文。如果输入过少，报告缺少依据；如果输入过多，模型容易遗漏重点或产生不稳定输出。

建筑幕墙韧性评估报告属于典型的多源信息整合任务。其输入包括项目基础信息、图像组信息、原子检测指标、图像产物、区域级数据、用户复核评论和任务状态。直接把这些信息一次性提交给大语言模型并不合理，因为底层指标数量可能随图像数量和区域数量快速增长。本文提出 LLM 分层信息整合策略，将报告生成拆分为图像级和项目级两个阶段。该策略并非简单压缩文本，而是根据工程评估逻辑先形成单图语义，再进行项目级融合。

\section{本章小结}
\label{sec:foundation-chapter-summary}

综上所述，现有研究在三个方面尚存在不足。第一，幕墙检测研究多关注单类缺陷识别，对图像资产管理、任务调度和报告生成关注不足。第二，Agent 研究多集中在通用工具调用、代码执行和问答任务，在建筑工程垂直场景中的系统化落地仍有探索空间。第三，大语言模型报告生成往往关注生成效果本身，而对输入信息如何分层组织、如何与数据库状态和工程流程结合讨论不足。

本文的研究切入点正是在上述空白之间建立连接：以幕墙检测为现实场景，以 Agent 作为任务调度和信息整合能力，以工程系统作为稳定运行载体，构建从图像资产到韧性评估报告的完整链路。

本文研究如何在真实业务流程中组织这些缺陷检测能力。已有检测模型回答“这张图是否存在某类缺陷”，本文系统进一步回答“这张图应该检测什么”“这些检测结果如何变成图像总结”“所有图像如何形成项目报告”。这三个问题正是 Agent 能力与工程链路结合的空间。
